\documentclass{article}

% Evaluations & Datasets Track, non-anonymous
\usepackage[eandd, nonanonymous]{neurips_2026}

\usepackage[utf8]{inputenc} % allow utf-8 input
\usepackage[T1]{fontenc}    % use 8-bit T1 fonts
\usepackage{hyperref}       % hyperlinks
\usepackage{url}            % simple URL typesetting
\usepackage{booktabs}       % professional-quality tables
\usepackage{amsfonts}       % blackboard math symbols
\usepackage{amsmath}        % \text in math
\usepackage{nicefrac}       % compact symbols for 1/2, etc.
\usepackage{microtype}      % microtypography
\usepackage{xcolor}         % colors
\usepackage{placeins}       % \FloatBarrier

% Unresolved markers from the draft ([QUOTE], [FILL], [CHECK], [X], ...)
\newcommand{\placeholder}[1]{\textcolor{red}{[#1]}}

% Drop the conference notice at the bottom of the first page
\makeatletter
\renewcommand{\@notice}{}
\makeatother

\title{GeoEnrich-NYC: A Geospatial Text-Tabular Dataset for Airbnb Price Prediction}

\author{%
  Author Name \\
  Affiliation \\
  \texttt{author@example.com} \\
}


\begin{document}


\maketitle


\begin{abstract}
  We provide a framework to enrich any dataset with existing latitude-longitude coordinates,
  where a regression task might benefit from a descriptive summary of each coordinates' surroundings.
\end{abstract}


\section{Introduction}

Accommodation datasets (Airbnb, booking.com) often include latitude-longitude coordinates
of listings.
The location of a listing can be leveraged as a significant indicator for price prediction
regression tasks, though it is unlikely that using the coordinates as-is will be an informative
feature for regression models.
To solve this problem, traditional ML regression models perform feature engineering and use
techniques such as HDBSCAN \citep{campello2013hdbscan} clustering, in order to leverage this potentially informative feature.

In addition, the surroundings of a listing can often provide signal regarding it's price. For example:
It is reasonable to assume \citep{wang2017price, zhang2017airbnb} that Airbnb listings in downtown Manhattan, where there is a high density
of restaurants tourist attractions, will be pricier than listings in the outskirts of the city -- where little
to no amenities can be found.

By following these assumptions, we hypothesize the following: a natural-language description of a listing's
surroundings will improve models' fit in a nightly-price prediction regression task.

To achieve this, we leverage the given latitude-longitude coordinates from an existing Airbnb listings dataset
and enrich them with geospatial data by querying Overture Maps places data \citep{overtureplaces} and enriching each listings with
business density information of it's surroundings. We then use batch LLM API calls \citep{openrouterbatch} to transform the
numerical density information into a description of each listing's surroundings.


\section{Related work}

\paragraph{Text-tabular curation and target-aware representations.} The goal of this project
is curating a dataset for the MulTaBench framework. MulTaBench frames text--tabular datasets
as candidates for curation rather than as evidence that any
single model is accurate: joint frozen inputs must improve on both unimodal
baselines, and target-aware representations (TAR) must improve on frozen ones
across several learners \citep{arazi2026multabench}.  Its text protocol begins with a
frozen pretrained encoder, such as E5 \citep{wang2022e5}, and adapts the encoder
to the prediction target with low-rank updates \citep[LoRA;][]{hu2022lora}.  We adopt
that protocol unchanged.  GeoEnrich-NYC contributes a new dataset candidate and
its reproducible curation pipeline; it neither proposes a tabular learner nor
claims that LLM verbalization is a replacement for target-aware representation
learning.

\paragraph{Geographic context in Airbnb pricing.} The basis of our approach is the idea that
geographic context provides a strong signal for nightly-price regression models.
In a personal project \citep{raz2026airbnbregression} I've used a similar dataset,
and performed feature engineering, transforming the raw coordinates into neighbourhoods,
derived using HDBSCAN clustering, and neighbourhood-median-price encoding as a feature.
This has proved to perform better than a city-wide encoding, which provides the base hypothesis
that a surroundings summary might provide similar value.
Similar approaches were applied in \citet{ahuja2021airbnb} and \citet{jiang2023multisource}, which achieved SOTA for this specific regression task,
and is generally common practice.

\paragraph{Prior Airbnb price studies} model
location directly through spatial weighting, coordinates, local amenities, or
spatial effects.  For example, \citet{zhang2017airbnb} use a geographically weighted
approach to study spatially varying price determinants, and \citet{tang2019spatial}
model spatial effects in peer-to-peer accommodation prices.  These studies support
the relevance of local context, but they do not test the MulTaBench question posed
here.  GeoEnrich-NYC instead uses coordinates only to construct evidence, removes
them from model inputs, and exposes a constrained textual account of nearby places
and landmarks.  This design tests a coarsened representation of location; it does
not estimate causal amenity effects or solve spatial dependence.

\paragraph{LLM enrichment.} Large language models are increasingly used as components of data pipelines,
not only as predictors. Previous studies have leveraged modern LLMs
both to structure unstructured or semi-structured data \citep{narayan2022wrangle} and to generate
a natural-language description of various features \citep{hegselmann2023tabllm}. These techniques
provide the basics upon which our solution is built.


\section{Dataset origins and curation methodology}

\subsection{Dataset and sources}

Inside Airbnb \citep{insideairbnb} is a project that provides quarterly data and advocacy
about airbnb's impact on residential communities, it's data is freely available
online. We used their New York City listings snapshot \citep{insideairbnbnyc}, scraped 14--23 June 2026, for the evaluation pipeline,
though the same pipeline can be applied to any other of their provided datasets.

Overture Maps Foundation is a non-profit organization that provides free and open-source maps data.
Overture publishes its data as GeoParquet files in a public S3 bucket (\texttt{s3://overturemaps-us-west-2}) \citep{overturedata}. We used DuckDB to query the
places theme of Overture release \texttt{2026-08-19.0} for businesses and points of interest in our
desired city, and performed a join by proximity for each listing.

\subsection{Pipeline}

\paragraph{Data cleaning.} After acquiring the original dataset from ``Inside Airbnb'', we perform a standard
feature engineering pipeline: columns that do not partake in the regression task are dropped,
the schema is defined and enforced, features such as ``number of rooms'' and ``numbers of baths''
are extracted from semi-structured text fields using simple regex rules. Listings with missing or non-positive prices are dropped,
as are listings missing coordinates or more than 25\% of their features.
In addition, the target nightly-price column of the final version of the dataset is log-transformed -- $\text{target} = \log(1 + price)$.
Of the 30,259 listings in the snapshot, 21,155 survive these filters; we randomly sample 15,000 of them (seed 0),
each with 8 feature columns (six numerical, one boolean, and one categorical) besides the target. The schema and feature distributions are summarized in Table~\ref{tab:data-card} (Appendix~\ref{app:data-card}).

\paragraph{Geospatial enrichment.} We query Overture for all places within the listings' bounding box (padded by $\approx 2$ km)
that have a confidence score of at least 0.6 and belong to one of our selected categories.
We apply a join such that each listing coordinate is assigned all POIs (points of interest) in a 450 meter radius, saved in JSON format.
The 10 listings with no POI within that radius are dropped, leaving 14,990 listings.
The original latitude-longitude columns are dropped.

\paragraph{Aggregation.} Each JSON columns describing nearby POIs is reduced down to a simple count,
and the distance to the nearest place, according to a list of 20 predefined buckets (for example: restaurants, parks, ground transport, nightlife venues, etc.)
Those categories are curated from the level-2 categories of the official Overture hierarchical taxonomy.
We rank each listing's count in every category against the citywide distribution of that count over all listings, and replace the count-buckets of each listing with it's relative percentile.
Popular POIs (landmarks) -- 77 curated attractions such as The Empire State Building, Statue of Liberty and Times Square -- are matched
within 800 meters of each listing and not reduced down to their categories, as they potentially carry stronger neighbourhood signal.

\paragraph{LLM descriptive summary.} By the end of the aggregation stage, each row has a JSON object containing the listing's
relative density of each category in their 450 meters surrounding radius and a list of popular tourist attractions nearby,
including their exact distance. We take a prompt blueprint (Appendix~\ref{app:prompt}) and structure the JSON summaries, preparing them for an LLM API call:
percentiles are phrased as word bands (e.g., ``far more than most neighbourhoods''), and only the three nearest landmarks are kept, with their distance phrased as a band (e.g., ``a short walk from'').
We send batch requests to an LLM, processing all 14,990 rows independently in three batches -- The LLMs are prompted to describe
the surroundings, and are not given the actual coordinates, exact counts, or listing price, to prevent data leakage.

Multiple LLMs were tested on subsets of the dataset, the model chosen is Gemini-3.5-flash-lite \citep{gemini35flashlite}. It was chosen for it's relative
lack of hallucinations, generation speed, batch requests availability and lack of reasoning -- as the price-per-row steeply increased for
reasoning models. A full model comparison is not in the scope of this project.

This pipeline results in 8 tabular feature columns and 1 textual column, describing a listing's surroundings.
Statistics regarding the text field can be found in Appendix~\ref{app:text-stats}.

\begin{figure}[h]
  \centering
  \small
  \fbox{\parbox{0.95\linewidth}{%
    \textbf{(1) Geospatial enrichment} --- POIs within 450\,m, as [count, nearest metres] per category (excerpt):\\[2pt]
    {\ttfamily\footnotesize
    restaurant: [83, 64] \quad casual\_eatery: [35, 54] \quad historic\_site: [24, 53] \quad discount\_store: [4, 324] \quad \ldots\\
    total places within 450\,m: 259 \quad landmarks within 800\,m: [Central Park, 792]}}}\\[2pt]
  $\downarrow$ \textit{citywide percentiles: total 259 $\rightarrow$ 60th; discount stores 4 $\rightarrow$ 97th; historic site 53\,m $\le$ 150\,m doorstep}\\[2pt]
  \fbox{\parbox{0.95\linewidth}{%
    \textbf{(2) Rendered evidence} --- the user message sent to the LLM (sections in per-listing shuffled order):\\[2pt]
    {\ttfamily\footnotesize
    Independent proximity evidence (use only if it adds a new idea):\\
    - historic sites: steps away\\[2pt]
    Overall: busy, by how many places are within a short walk compared with other NYC neighbourhoods.\\[2pt]
    Primary evidence --- strongest citywide contrasts:\\
    - discount stores: far more than most neighbourhoods\\[2pt]
    Well-known places nearby (copy names exactly, do not add any):\\
    - about 10 minutes from: Central Park}}}\\[2pt]
  $\downarrow$ \textit{Gemini-3.5-flash-lite, prompt 16 system instruction (Appendix~\ref{app:prompt})}\\[2pt]
  \fbox{\parbox{0.95\linewidth}{%
    \textbf{(3) \texttt{surroundings\_summary}:}
    ``This busy neighbourhood features far more discount stores than most neighbourhoods, with historic sites steps away.
    Central Park is about 10 minutes away.''}}
  \caption{How a row turns into a description, for one listing of the released dataset (\texttt{index} 10433; entire home, 7 guests).
    The coordinates, exact counts, distances and price never reach the LLM; only the word bands in (2) do.}
  \label{fig:row-to-description}
\end{figure}


The process involved trial and error with numerous prompt versions,
tested on a small subset ($\approx 2000$ rows) of the dataset, until we settled on the 16th variant.
Once the final variant has been chosen, we have added surrounding descriptions to the entire 14,990 rows.
We then ran the official MulTaBench \citep{arazi2026multabench} curation pipeline once, with log-transformed price
as the target column.


\section{Results}

We evaluate GeoEnrich-NYC with the official MulTaBench benchmark, run unmodified, with the
log-transformed nightly price $\log(1 + \text{price})$ as the target.
Each of the five committee learners (LightGBM, CatBoost, TabM, TabPFNv2, TabPFNv2.5) is trained under
four conditions: structured columns only, the surroundings description only (text-only),
structured columns with frozen E5-small \citep{wang2022e5} text embeddings (joint frozen), and structured columns with
E5-small embeddings fine-tuned on the target with LoRA \citep{hu2022lora} (joint TAR).
Every condition runs on the benchmark's six splits of TabArena's 67/33 scheme
($\approx 9{,}993$ training and $\approx 4{,}997$ test rows per split).
The metric is the RMSE on the log-price scale; lower is better.

A learner shows \emph{joint signal} when its mean joint-frozen RMSE is lower than both its structured and
its text-only RMSE by more than $0.001$, and \emph{TAR gain} when its mean joint-TAR RMSE is lower than its
joint-frozen RMSE by more than $0.001$ (means rounded to three decimals before differencing).
The dataset is eligible when at least 3 of the 5 learners show both.

\begin{table}[h]
  \caption{RMSE on $\log(1 + \text{price})$, mean $\pm$ standard deviation over six splits (lower is better).
    $\Delta_{\text{joint}}$: improvement of joint frozen over the better unimodal baseline.
    $\Delta_{\text{TAR}}$: improvement of joint TAR over joint frozen.
    All five learners show both joint signal and TAR gain (5/5); the dataset is eligible.}
  \label{tab:results}
  \centering
  \small
  \begin{tabular}{lccccrr}
    \toprule
    Learner & Structured & Text-only & Joint frozen & Joint TAR & $\Delta_{\text{joint}}$ & $\Delta_{\text{TAR}}$ \\
    \midrule
    LightGBM   & 0.501 $\pm$ 0.007 & 0.726 $\pm$ 0.006 & 0.442 $\pm$ 0.006 & 0.435 $\pm$ 0.005 & +0.059 & +0.007 \\
    CatBoost   & 0.498 $\pm$ 0.006 & 0.722 $\pm$ 0.008 & 0.435 $\pm$ 0.005 & 0.431 $\pm$ 0.007 & +0.063 & +0.004 \\
    TabM       & 0.505 $\pm$ 0.008 & 0.720 $\pm$ 0.009 & 0.434 $\pm$ 0.004 & 0.432 $\pm$ 0.005 & +0.071 & +0.002 \\
    TabPFNv2   & 0.490 $\pm$ 0.008 & 0.715 $\pm$ 0.007 & 0.422 $\pm$ 0.005 & 0.418 $\pm$ 0.006 & +0.068 & +0.004 \\
    TabPFNv2.5 & 0.490 $\pm$ 0.007 & 0.710 $\pm$ 0.008 & 0.415 $\pm$ 0.004 & 0.412 $\pm$ 0.005 & +0.075 & +0.003 \\
    \bottomrule
  \end{tabular}
\end{table}

Table~\ref{tab:results} reports the results.
Adding the surroundings description to the structured columns lowers RMSE for every learner,
by $0.059$--$0.075$ ($11.7$--$15.3\%$ relative to the structured baseline), so all five learners show joint signal.
The description alone is a weak predictor: text-only RMSE ($0.710$--$0.726$) is well above the structured baseline
for every learner, which suggests the text complements the listing attributes rather than substituting for them.
Fine-tuning the encoder on the target lowers RMSE further for all five learners, so all five show TAR gain and the
dataset passes the MulTaBench curation criterion with 5/5 learners on both checks.
The TAR gain is small, however ($0.002$--$0.007$), and of the same order as the split-to-split standard deviation;
joint TAR beats joint frozen on 27 of the 30 learner--split pairs (6/6 splits for LightGBM, TabPFNv2 and TabPFNv2.5;
5/6 for CatBoost; 4/6 for TabM).
The best configuration, TabPFNv2.5 with joint TAR, reaches an RMSE of $0.412$, a $16.1\%$ reduction over its structured baseline.


\section{Conclusion, limitations, and future work}

\subsection{Outcomes}

We have successfully passed the MulTaBench curation pipeline, passing 5/5 learners on joint signal,
and 5/5 learners on TAR signal.
The constructed framework is flexible, and could potentially enrich any number of datasets, provided that:
coordinates are already provided, and the regression / classification task could potentially benefit from a
description of the coordinates surroundings.

\subsection{Limitations}

\paragraph{Evaluation.} Though the evaluation pipeline passes successfully, the reason behind that could be further explored.
It is unclear what the main driver of the price-signal is in each listing's surroundings description.
It is possible that it is derived by the explicit naming of popular landmarks (present in the evidence of 42.6\% of listings),
or perhaps by the overall busyness band and the category contrasts.
The pipeline already builds two censored arms that separate these sources --- one keeping only the landmark names found in each summary,
and one replacing each landmark name with ``a well-known landmark'' --- but they were not evaluated in this release.

\paragraph{Application.} Though the provided framework is flexible and could potentially be applied to any number of datasets,
the structure of the prompts make it more applicable for hospitality or real-estate related datasets.
In addition, the landmarks were curated manually and so applying the pipeline to a different city will require
updating that list, or perhaps automating that choice through some different API call.

\subsection{Future work}

This project only tackles a single city, while the pipeline can potentially be applied across all datasets with regression tasks
that could benefit from geospatial information. Future work could evaluate the robustness of this pipeline across datasets outside of
hospitality sector price regression tasks.

A major challenge of this project was increasing the variance of responses (measured as the share of distinct summaries
and the mean pairwise Jaccard similarity of their content words; Appendix~\ref{app:text-stats}) while maintaining the TAR gain.
The released corpus has 95.1\% distinct summaries but a vocabulary of only 442 word types, so many listings are still
described with near-identical phrasing.
Future work could tackle increasing variance, by using several prompt blueprints, randomly assigning dataset rows to different LLMs, or
using more powerful models.


{\small
\bibliographystyle{plainnat}
\bibliography{references}
}


%%%%%%%%%%%%%%%%%%%%%%%%%%%%%%%%%%%%%%%%%%%%%%%%%%%%%%%%%%%%

\appendix

\section{Data card}
\label{app:data-card}

\begin{table}[h]
  \caption{Schema and feature distributions of the released dataset (14,990 rows). Quantiles are over non-missing values.
    The target is $\log(1+\text{price})$, with price the nightly price in USD (raw quantiles: 7, 96, 173, 299, 30{,}973).}
  \label{tab:data-card}
  \centering
  \small
  \setlength{\tabcolsep}{3pt}
  \begin{tabular}{llrrrrrr}
    \toprule
    Column & Type & Missing & Min & Q1 & Median & Q3 & Max \\
    \midrule
    \texttt{price} (target) & log USD & 0.0\% & 2.10 & 4.58 & 5.16 & 5.70 & 10.34 \\
    \texttt{ratings} & numeric & 27.8\% & 1.00 & 4.67 & 4.86 & 5.00 & 5.00 \\
    \texttt{guests} & numeric & 0.0\% & 1 & 2 & 2 & 4 & 16 \\
    \texttt{bathrooms} & numeric & 3.6\% & 0 & 1 & 1 & 1 & 15.5 \\
    \texttt{num\_bedrooms} & numeric & 33.8\% & 0 & 1 & 1 & 2 & 19 \\
    \texttt{property\_num.\_of\_reviews} & numeric & 0.0\% & 0 & 0 & 6 & 41 & 4502 \\
    \texttt{minimum\_nights} & numeric & 0.0\% & 1 & 30 & 30 & 30 & 365 \\
    \midrule
    \texttt{room\_type} & categorical & 0.0\% & \multicolumn{5}{l}{\begin{tabular}[t]{@{}l@{}}Entire home 54.7\%, Private 42.1\%,\\ Hotel 2.3\%, Shared 0.9\%\end{tabular}} \\
    \texttt{is\_superhost} & boolean & 0.0\% & \multicolumn{5}{l}{True 31.1\%} \\
    \texttt{surroundings\_summary} & text & 0.0\% & \multicolumn{5}{l}{14,254 distinct} \\
    \bottomrule
  \end{tabular}
\end{table}

\begin{table}[h]
  \caption{Row attrition from the Inside Airbnb snapshot to the released dataset, and duplicates.}
  \label{tab:attrition}
  \centering
  \small
  \begin{tabular}{lr}
    \toprule
    Step & Rows \\
    \midrule
    NYC snapshot (scraped 14--23 June 2026) & 30,259 \\
    Positive, non-missing price & 21,515 \\
    Coordinates present and at most 25\% of features missing & 21,155 \\
    Random sample (seed 0) & 15,000 \\
    At least one POI within 450\,m (released) & 14,990 \\
    \midrule
    Rows whose 8 tabular features duplicate another row & 6,066 \\
    Rows whose 8 tabular features and summary duplicate another row & 46 \\
    \bottomrule
  \end{tabular}
\end{table}

\section{Prompt blueprint}
\label{app:prompt}

Each request has the system instruction below and a user message rendered from the listing's aggregated surroundings
(Figure~\ref{fig:row-to-description}, panel 2). The user message has four sections, shown in an order shuffled per listing with
a seed derived from its surroundings, so it is reproducible: the overall busyness, up to four primary contrasts
(at most one per Overture level-1 group), at most one doorstep proximity cue, and up to three landmarks.

\begin{quote}
  \small
  \textbf{System instruction (prompt 16).}
  Write one or two neutral, factual sentences describing the surroundings of an Airbnb listing in New York City, using only the evidence given.
  The ``Overall'' line sets the tone: a quiet or very quiet area must read as quiet, never as busy, lively, walkable, convenient, or full of amenities; a busy area may read as busy.
  State scarcity plainly (``few restaurants'', ``no gyms nearby''); do not dramatise it (``remarkably isolated'', ``peaceful retreat'') and do not infer anything not listed, such as a residential, commercial, or nightlife character.
  A proximity cue means one such place is close by, not that the area has many of them.
  Keep the strength of each given comparison (``more than most'', ``far more than most'', ``more than almost anywhere''); do not intensify or soften it, and add no strong words of your own.
  Plain text, no markdown.
\end{quote}

\begin{table}[h]
  \caption{Word bands used to render the evidence. Percentiles are mid-rank (ties counted half) against all 14,990 listings.
    A category is a primary contrast only in the outer 15\% on either side; a high band needs at least 3 places,
    and ``none nearby'' is used only for categories present near at least 60\% of listings.}
  \label{tab:bands}
  \centering
  \small
  \begin{tabular}{p{0.17\linewidth}p{0.24\linewidth}p{0.49\linewidth}}
    \toprule
    Evidence & Input & Band \\
    \midrule
    Overall busyness & percentile of all places within 450\,m
      & $<$20: very quiet; $<$40: quiet; $<$60: average; $<$80: busy; else very busy \\
    Primary contrast (high) & category percentile $\ge 85$
      & $\ge$99: more than almost anywhere in NYC; $\ge$95: far more than most; else more than most \\
    Primary contrast (low) & category percentile $\le 15$
      & 0 places: none nearby; $\le$5: very few; else fewer than most \\
    Proximity cue & nearest place $\le 150$\,m
      & $\le$50\,m: on the doorstep; $\le$150\,m: steps away \\
    Landmark & distance $\le 800$\,m
      & $\le$100: right by; $\le$250: a couple minutes from; $\le$500: a short walk from; else about 10 minutes from \\
    \bottomrule
  \end{tabular}
\end{table}

Proximity cues are limited to parks, ground transport, museums, hotels, nightlife venues and historic sites,
and skip any category already listed as a primary contrast.

\begin{table}[h]
  \caption{Generation settings and cost of the released corpus.}
  \label{tab:generation}
  \centering
  \small
  \begin{tabular}{ll}
    \toprule
    Setting & Value \\
    \midrule
    Model & \texttt{gemini-3.5-flash-lite-20260721} (OpenRouter batch API) \\
    Temperature / max output tokens & 0.4 / 300 \\
    Batches & 3 (5,000 + 5,000 + 4,990 requests); 0 failed \\
    Tokens & 4,335,957 prompt + 686,500 completion \\
    Cost & \$1.51 total \\
    \bottomrule
  \end{tabular}
\end{table}

\FloatBarrier
\section{Summary text statistics}
\label{app:text-stats}

\begin{table}[h]
  \caption{Statistics of \texttt{surroundings\_summary} over the 14,990 released rows. Words are alphabetic tokens.
    Jaccard similarity is computed on content-word sets (stopwords and generic words such as ``area'' and ``neighborhood'' removed)
    over 20,000 random pairs (seed 0); lower means more varied text.}
  \label{tab:text-stats}
  \centering
  \small
  \begin{tabular}{lr}
    \toprule
    Statistic & Value \\
    \midrule
    Words per summary, mean $\pm$ SD (median; min--max) & 39.1 $\pm$ 11.1 (38; 11--84) \\
    Characters per summary, mean (median; max) & 238.7 (229; 522) \\
    Sentences per summary: 1 / 2 / 3 & 450 / 14,071 / 469 \\
    Distinct summaries & 14,254 (95.1\%) \\
    Distinct rendered evidence blocks & 12,594 (84.0\%) \\
    Vocabulary (lower-cased word types) & 442 \\
    Mean pairwise content-word Jaccard similarity & 0.185 \\
    Listings with a landmark in the evidence & 42.6\% \\
    \quad of which the summary names at least one landmark & 99.6\% \\
    Summaries containing ``quiet'' / ``busy'' & 39.4\% / 39.9\% \\
    \bottomrule
  \end{tabular}
\end{table}

\FloatBarrier
\section{Tables and metrics only}
\label{app:tables}

Do not add prose, qualitative examples, figures, or new claims. Every appendix table must support a main-paper claim or artifact audit.

\begin{table}[h]
  \caption{Appendix tables and their required content.}
  \label{tab:appendix-plan}
  \centering
  \begin{tabular}{lp{0.8\linewidth}}
    \toprule
    Table & Required content \\
    \midrule
    A2 & Source/provenance/reuse matrix: URL, snapshot, retrieval date, terms/license, transformations, redistribution status \\
    A3 & Reproducibility matrix: folds, seeds, software, hyperparameters, PCA/encoder/LoRA, runtime, hardware, compute/environmental estimate \\
    A4 & Full per-fold primary results, deltas, uncertainty, and pass flags \\
    A5 & Raw-aggregate, genuine-text, mismatched-text, named-landmark, and redacted-landmark diagnostic results \\
    A6 & Spatial-holdout construction, sample counts, and results \\
    A7 & Grounding/privacy audit counts, named-versus-redacted re-identification checks, and artifact manifest \\
    \bottomrule
  \end{tabular}
\end{table}

%%%%%%%%%%%%%%%%%%%%%%%%%%%%%%%%%%%%%%%%%%%%%%%%%%%%%%%%%%%%

\end{document}
